% 第8章“异常结构检测与修复多场景案例验证”完整替换稿
% 数据来源：14个BPMN协作模型和221个日志输入，分别用于案例有效性、规则覆盖和规模性能验证。
% 建议替换原文中从本节标题开始，到“与现有偏差检测、模型修复方法的对比”小节之前的内容。

\section{异常结构检测与修复多场景案例验证}
\label{sec:异常结构检测与修复多场景案例验证}

本节以已经通过协作业务流程验证的BPMN协作图为输入，对全局消息异常检测、最小异常结构定位和潜在修复方案生成方法进行多场景案例验证。与只报告最终判定的实验方式不同，本次实验同时记录实际输入文件、消息迹投影、消息模式判定、预分割过程、最小异常结构（Minimal Abnormal Structure，MinAS）、修复规则来源以及修复前后的消息迹复验结果。实验围绕以下问题展开：

\begin{itemize}[nosep]
    \item 研究问题1：本文方法能否根据全局消息日志锁定潜在异常子流程，并进一步定位MinAS？
    \item 研究问题2：本文方法能否在不同参与方数量、消息流数量和控制流结构的协作场景中稳定执行？
    \item 研究问题3：当协作流程规模和并行结构复杂度增加时，完整检测与修复分析流水线的耗时如何变化？
    \item 研究问题4：潜在修复方案能否在接纳异常消息迹的同时保留原模型行为和全部已观测正常消息迹？
\end{itemize}

\subsection{实验设计与复现设置}

实验数据集包含14个不同的BPMN协作模型和221个全局消息日志输入，并根据验证目的划分为多场景人工案例组、业务场景规则覆盖组和规模化合成组。多场景人工案例组包含8个BPMN和67个日志，覆盖青岛港出口、混凝土供应链合规检测、机场护照检查、供应链付款发货、自助餐厅、订单报价和银行信用评估等场景。其中，青岛港出口协作图（图\ref{fig:港口协作图（业务约束）}）来源于与实际业务设计人员的多次讨论；其余模型来自相关工作、Camunda和BPMAI公开案例\cite{spalazzi2023blockchain,compagnucci2024study,camunda_bpmn_repository,bpmai_repository}，并按照本文在第\ref{sec:协作业务流程与业务约束建模}节给出的结构化、无循环和BPMN元素子集约束进行了适配。业务场景规则覆盖组在上述8个BPMN上配置64个受控变异日志；规模化合成组包含6个结构化BPMN和90个日志，用于考察参与方、消息流及并行结构增加后的性能变化。

每次实验以一个BPMN文件和一个日志文件为完整输入。日志文件包含若干正常消息迹以及人工构造或程序注入的异常消息迹。程序既支持运行时从\path{experiments}目录选择案例、BPMN文件和日志文件，也支持通过命令行参数固定输入。为保证结果可追溯，程序在运行开始时输出两个输入文件的绝对路径、文件大小和SHA-256摘要；在运行结束后，将逐参与方投影、消息模式判定、预分割步骤、MinAS、修复候选和复验统计写入\path{analysis-report.json}。本次批量实验使用默认参数\texttt{pattern-limit=10000}和\texttt{behavior-limit=20000}。其中，前者只限制消息模式的展示数量，不影响精确成员判定；后者限制定义\ref{def:潜在修复方案}所需的子树行为枚举，达到上限时程序明确返回“无法证明”，而不将其当作修复成功。

实验于Windows 10（10.0.19045）、AMD Ryzen 5 5600处理器、Python 3.12.4和PM4Py 2.7.19.8环境中执行。每个日志文件均通过一次独立程序运行完成分析。表\ref{tab:多场景协作业务流程异常检测与修复案例实验结果}中的耗时直接取JSON报告的\texttt{elapsed\_seconds}字段，并计算同一场景全部日志输入的算术平均值。该字段从分析函数入口开始计时，包括输入读取、BPMN规范化、流程树转换、消息模式计算、异常定位、修复候选生成、复验和输入摘要计算，但不包括Python模块导入、JSON序列化写入和终端报告打印。

为避免把“存在一个可接纳异常的改动”直接等同于“修复有效”，本节采用以下严格判定标准。

\begin{definition}
\label{def:异常修复方案严格有效性}
异常修复方案严格有效性：对于一个修复作用域，若至少存在一个候选方案同时满足：（1）修复前消息模式是修复后消息模式的子集，即保留原模型行为；（2）覆盖该作用域中的全部异常观测；（3）将候选应用到完整子流程后，能够重放该输入中的全部正常消息迹和全部异常消息迹，则称该作用域存在严格有效修复候选。只有当一次日志输入产生的每个修复作用域都存在上述候选时，才将该日志输入记为“整体修复成功”。
\end{definition}

表\ref{tab:多场景协作业务流程异常检测与修复案例实验结果}中的“成功锁定”表示该日志输入至少识别出一个潜在异常子流程；该指标反映当前程序对人工异常输入的可观测检出情况，不等同于带人工位置标注数据集上的定位准确率。“严格修复成功”采用定义\ref{def:异常修复方案严格有效性}，因而比仅要求修复后通过异常消息迹的判定更严格。

\begin{table}[htbp]
\centering
\small
\renewcommand{\arraystretch}{1.25}
\renewcommand{\tabularxcolumn}[1]{m{#1}}
\caption{多场景人工案例组的异常检测与修复实验结果}
\label{tab:多场景协作业务流程异常检测与修复案例实验结果}
\begin{tabularx}{\textwidth}{
    >{\centering\arraybackslash\hsize=1.55\hsize}X
    >{\centering\arraybackslash\hsize=0.85\hsize}X
    >{\centering\arraybackslash\hsize=0.70\hsize}X
    >{\centering\arraybackslash\hsize=0.95\hsize}X
    >{\centering\arraybackslash\hsize=0.95\hsize}X
    >{\centering\arraybackslash\hsize=0.70\hsize}X
}
\toprule
BPMN协作业务流程 & 子流程数、消息流数 & 日志输入数 & 成功锁定潜在异常子流程 & 严格修复成功 & 平均分析耗时（秒） \\
\midrule
青岛港出口协作图（图\ref{fig:港口协作图（业务约束）}） & 9，25 & 27 & 27/27 & 16/27 & 4.1403 \\
青岛港出口简化协作图（图\ref{fig:港口协作图简化}） & 3，8 & 7 & 6/7 & 5/7 & 3.9890 \\
混凝土供应链合规检测协作流程\cite{spalazzi2023blockchain} & 4，10 & 7 & 7/7 & 3/7 & 4.0113 \\
机场护照检查流程\cite{bpmai_repository} & 2，6 & 6 & 6/6 & 4/6 & 3.9671 \\
供应链付款发货协作流程\cite{compagnucci2024study} & 2，7 & 5 & 4/5 & 4/5 & 3.9604 \\
自助餐厅业务流程\cite{camunda_bpmn_repository} & 3，7 & 6 & 6/6 & 6/6 & 3.9708 \\
供应链订单报价协作流程\cite{camunda_bpmn_repository} & 2，4 & 4 & 4/4 & 2/4 & 3.9617 \\
银行信用评估流程\cite{camunda_bpmn_repository} & 3，6 & 5 & 5/5 & 4/5 & 3.9885 \\
\midrule
合计或加权平均 & -- & 67 & 65/67 & 44/67 & 4.0449 \\
\bottomrule
\end{tabularx}
\end{table}

\subsection{实验数据集构成与受控生成方法}
\label{subsec:实验数据集构成与受控生成方法}

如表\ref{tab:实验数据集构成}所示，多场景人工案例组使用67个包含人工异常的日志评价方法在不同业务场景中的可观测检出能力和严格修复能力；业务场景规则覆盖组使用相同的8个BPMN，通过对能够被全部参与方接受的正常消息迹执行消息删除、相邻消息换序以及发送端/接收端可见性变异，形成64个受控日志；规模化合成组生成6个结构化、无循环的BPMN协作图，参与方数量由5增加到20，消息流数量由12增加到96，并为每个模型配置15个日志，共90个输入。三组数据合计包含14个不同BPMN和221个日志输入。

受控数据由脚本\path{scripts/generate_extended_benchmark.py}在固定随机种子20260809下生成。对于规模化合成模型，协调者子流程由顺序块、并行块和排他选择块组合而成，其余参与方通过消息流与协调者交互。每个模型分别配置5个消息缺失异常、5个顺序逆转异常和5个互斥分支同时执行异常。对于业务场景模型，生成器首先从人工案例日志中筛选被全部参与方接受的正常迹，再对其执行候选变异。每个受控日志都包含生成种子、异常类型和注入位置注释；JSON和CSV清单记录BPMN及日志的SHA-256、受影响消息、检出的参与方、MinAS表达式和通过严格复验的表7-5规则编号。

\begin{table}[htbp]
\centering
\small
\renewcommand{\arraystretch}{1.25}
\caption{实验数据集构成及其验证用途}
\label{tab:实验数据集构成}
\begin{tabularx}{\textwidth}{
    >{\centering\arraybackslash\hsize=1.25\hsize}X
    >{\centering\arraybackslash\hsize=0.65\hsize}X
    >{\centering\arraybackslash\hsize=0.75\hsize}X
    >{\centering\arraybackslash\hsize=0.80\hsize}X
    >{\centering\arraybackslash\hsize=1.55\hsize}X
}
\toprule
数据组 & BPMN数量 & 日志输入数 & 模型规模范围（参与方、消息流） & 主要用途 \\
\midrule
多场景人工案例组 & 8 & 67 & $(2,4)$--$(9,25)$ & 场景有效性评价和局限性分析 \\
业务场景规则覆盖组 & 8（与案例组相同） & 64 & $(2,4)$--$(9,25)$ & 修复规则覆盖和回归稳定性验证 \\
规模化合成组 & 6 & 90 & $(5,12)$--$(20,96)$ & 受控正确性验证和规模性能分析 \\
\midrule
合计 & 14个不同BPMN & 221 & $(2,4)$--$(20,96)$ & -- \\
\bottomrule
\end{tabularx}
\end{table}

154个受控日志的异常类型分布为：消息缺失51个、顺序逆转30个、互斥分支同时执行30个、业务场景相邻消息换序9个、仅发送端可见16个、仅接收端可见18个。独立验证脚本\path{scripts/validate_extended_benchmark.py}重新加载每个BPMN，校验文件摘要和消息迹数量，并重新计算异常参与方、MinAS和严格有效规则。154个受控输入均通过独立复核。此外，对6个规模化合成模型的3类异常各抽取1个输入，并对8个业务场景模型各抽取1个输入，共执行26次完整分析流水线冒烟测试，未发现结果不一致。

需要指出的是，154个受控输入按照本文方法能够表达的异常类型设计，并在纳入数据集前经过严格有效性筛选。因此，受控数据的154/154检出和修复结果用于证明实现对目标规则的覆盖能力、验证回归稳定性并开展规模实验，不能被解释为未知分布上的无偏准确率。对于实际场景中的有效性问题，本文采用多场景人工案例组的$65/67$可观测检出比例和$44/67$严格日志级修复比例；对于规则覆盖和性能问题，则分别采用业务场景规则覆盖组和规模化合成组，以避免由样本筛选造成评价偏差。

\subsection{规模性能实验}
\label{subsec:规模性能实验}

规模性能实验使用\path{scripts/benchmark_extended_scale.py}完成。在正式计时前，先执行一次PM4Py模型加载作为预热；随后对6个合成BPMN的全部90个日志各执行一次完整分析。计时直接采用程序生成的\texttt{elapsed\_seconds}字段，包括BPMN规范化、流程树转换、消息模式分析、MinAS求解、修复候选生成、严格复验和输入摘要计算，不包括Python及PM4Py模块导入、JSON写入和终端输出。实验环境为Windows 10、AMD Ryzen 5 5600、Python 3.12.4和PM4Py 2.7.19.8。每个规模包含15次运行，分别覆盖5个消息缺失、5个顺序逆转和5个选择冲突输入。

\begin{table}[htbp]
\centering
\small
\renewcommand{\arraystretch}{1.25}
\caption{规模递增合成BPMN的预热分析性能}
\label{tab:规模递增合成BPMN性能}
\begin{tabularx}{\textwidth}{
    >{\centering\arraybackslash}X
    >{\centering\arraybackslash}X
    >{\centering\arraybackslash}X
    >{\centering\arraybackslash}X
    >{\centering\arraybackslash}X
    >{\centering\arraybackslash}X
    >{\centering\arraybackslash}X
}
\toprule
规模编号 & 参与方数 & 消息流数 & 日志数 & 检出/严格修复 & 平均耗时（秒） & P95（秒） \\
\midrule
Scale-1 & 5 & 12 & 15 & 15/15 & 0.0768 & 0.1061 \\
Scale-2 & 7 & 20 & 15 & 15/15 & 0.1252 & 0.1422 \\
Scale-3 & 9 & 32 & 15 & 15/15 & 0.2734 & 0.3560 \\
Scale-4 & 12 & 48 & 15 & 15/15 & 0.5806 & 0.6967 \\
Scale-5 & 16 & 72 & 15 & 15/15 & 2.3086 & 2.4601 \\
Scale-6 & 20 & 96 & 15 & 15/15 & 5.1961 & 5.4534 \\
\bottomrule
\end{tabularx}
\end{table}

如表\ref{tab:规模递增合成BPMN性能}所示，在参与方数量由5增加至20、消息流数量由12增加至96时，平均分析耗时由0.0768秒增加至5.1961秒。前4个规模的平均耗时均低于0.6秒；当消息流增加至72和96时，耗时分别上升至2.3086秒和5.1961秒。这一变化表明，当前实现在中小规模模型上开销较低，但包含更多并行块的较大模型会显著增加消息模式展示枚举、流程树转换和修复行为验证的计算量。最大模型的15次运行P95为5.4534秒，且15个输入均完成异常定位和严格修复复验，没有发生超时或枚举结果被静默截断。

从Scale-1到Scale-6，消息流数量增加8倍，而平均耗时约增加67.6倍，说明实验区间内的运行时间并非随消息流数量线性增长。该结果揭示了并行组合空间增加所带来的性能代价，但仍不能直接作为算法渐近复杂度的证明。未来可以通过固定控制流结构、分别改变参与方数量、消息流数量和并行块数量的消融实验，进一步区分各因素对耗时和内存占用的影响。




\subsection{多场景人工案例实验结果及分析}

\subsubsection{潜在异常子流程锁定与MinAS定位}

67个日志输入均完成了BPMN解析和三阶段分析，没有发生程序运行失败。其中，65个输入至少锁定了一个潜在异常子流程，对人工标记异常输入的可观测检出比例为$65/67=97.0\%$。两个未锁定输入分别为青岛港简化案例的\texttt{global\_log\_6.txt}和供应链付款发货案例的\texttt{global\_log\_2.txt}。前者的异常迹仅将最后一条消息记为\texttt{M8\_s}；按发送端/接收端方向投影后，货代、船代和港口三个参与方分别得到$(M1,M2,M3,M4)$、$(M1,M5,M2,M3,M6,M7)$和$(M4,M6,M5,M7,M8)$，均属于相应的open或closed消息模式。后者的异常迹为$(Order,Notification\_Nok\_s)$；投影后，客户方只观察到$(Order)$，公司方观察到$(Order,Notification\_Nok)$，两者同样被现有消息模式接受。因此，这两个结果不是运行错误，而是表明仅凭当前方向化消息投影和open/closed消息模式不能区分这两类观测。该反例说明，方法的适用性受公共消息可观测性和消息模式抽象粒度约束。

对于成功锁定的输入，程序先将全局消息迹投影到每个参与方的消息集合，再依据表\ref{tab:消息模式求解规则}进行精确成员判定；对不属于消息模式的投影，程序按照表\ref{tab:预分割规则}和算法\ref{alg:从异常子流程中检测最小异常结构}逐层执行预分割，并记录节点路径、当前子树表达式、分配到各子节点的消息片段、匹配结果和失败原因。与旧实现只返回一个活动名称相比，新实现能够把异常定位到叶子节点、完整控制流子树或顺序节点的连续子区间，并按照定义\ref{def:合并最小异常结构}合并相关MinAS。

以青岛港简化案例\texttt{global\_log\_1.txt}为例，程序实际使用\path{experiments/qingdao_port_simple/collaboration.bpmn}和\path{experiments/qingdao_port_simple/global_log_1.txt}。货代子流程的消息流程树为$.(M1,M2,M3,M4)$，异常迹投影为$(M1,M3,M2,M4)$。预分割在$M2$和$M3$所在的连续子区间发现顺序逆转，因而把MinAS定位为$.(M2,M3)$，而不是把整个货代子流程误判为异常结构。船代子流程也得到相同的局部MinAS，港口子流程则保持正常。

\begin{lstlisting}[caption={新原型系统对顺序逆转异常的输出节选},label={lst:detect_repair_new}]
[输入]
 BPMN: experiments/qingdao_port_simple/collaboration.bpmn
 日志: experiments/qingdao_port_simple/global_log_1.txt
 （程序同时输出绝对路径、大小和 SHA-256）

[参与方 forwarder]
 流程树: .(M1,M2,M3,M4)
 异常投影: (M1,M3,M2,M4)
 预分割结果: M2 与 M3 的顺序约束不满足
 最小异常结构: .(M2,M3)

[候选1] 表7-5规则5
 .(M2,M3) -> |(M2,M3)
 定义7.17: 满足
 正常消息迹复验: 7/7
 异常消息迹复验: 1/1

[候选2] 表7-5规则4
 .(M2,M3) -> .(M3,M2)
 定义7.17: 不满足
 正常消息迹复验: 0/7
 异常消息迹复验: 1/1
\end{lstlisting}

该结果还展示了严格复验的必要性。规则4虽然能够接纳当前逆序异常迹，但会使7条正常迹全部失效；规则5将局部顺序结构改为并行结构，既保留正序行为，也接纳逆序行为，因此被排在规则4之前。另一个代表性输入\texttt{global\_log\_5.txt}中，港口子流程异常投影为$(M4,M6,M5,M8)$，系统将缺失的$M7$定位为叶子MinAS，并依据表7-5规则2生成$M7\rightarrow +(M7,\tau)$。该候选通过7/7条正常迹和1/1条异常迹；直接删除$M7$的规则1候选只能通过异常迹而不能通过正常迹，因此不满足严格有效性。

\subsubsection{修复有效性与规则分布}

67个日志输入共生成118个修复作用域，其中94个作用域至少存在一个严格有效候选，作用域级成功比例为$94/118=79.7\%$。若按一次日志输入的所有修复作用域整体计算，则44个输入全部通过，日志级严格修复成功比例为$44/67=65.7\%$；若只以已经锁定潜在异常子流程的65个输入为分母，则为$44/65=67.7\%$。两种比例分别反映完整实验输入和已检出输入上的修复能力，本文主要采用前者，以避免排除未检出案例造成选择偏差。

通过严格复验的候选中，表7-5规则2“消息活动改为可选”出现84次，规则5“顺序改为并行”出现7次，规则8“排他选择改为并行”出现3次。其余候选可能能够覆盖当前异常片段，但由于删除了原有行为、不能通过全部正常迹，或者只解决了同一日志中的部分修复作用域，未计为严格有效。该结果表明，潜在修复方案的价值不仅在于给出一个能够拟合异常迹的模型改动，还在于把“为何推荐”和“为何淘汰”都转化为可检查的行为包含关系与日志复验证据。与此同时，44/67的日志级结果也说明，表7-5的局部单作用域规则尚不能组合解决所有复合异常，后续工作需要研究多MinAS候选的联合应用、冲突消解和全局排序。

\subsubsection{执行时间}

多场景人工案例组67次独立运行的分析流水线平均耗时为4.0449秒，标准差为0.1841秒，最小值为3.9010秒，最大值为5.2563秒。规模最大的青岛港完整模型平均耗时为4.1403秒，其余场景均值位于3.9604秒至4.0113秒之间。各场景耗时较为接近，说明BPMN规范化和流程树转换等固定开销占有较大比例；仅凭这一组模型不能推导更大规模下的性能。第\ref{subsec:规模性能实验}节使用6个规模化合成BPMN和90个日志执行预热性能实验，专门分析参与方、消息流和并行组合空间增加后的耗时变化。

\subsubsection{研究问题回答}

对于研究问题1，本文方法在多场景人工案例组的65/67个异常日志输入上锁定了潜在异常子流程，并能够给出节点路径、MinAS表达式和预分割证据；两个方向化消息观测反例表明，检测能力受消息可观测性和open/closed模式语义限制。对于研究问题2，14个BPMN和221个日志输入均能完成分析，说明原型系统可以处理不同数量的参与方、消息流及控制流结构；其中154个受控输入用于验证异常类型和修复规则覆盖，不能与人工案例组的有效性比例混合解释。对于研究问题3，多场景人工案例组的独立进程分析流水线平均耗时约为4.04秒；规模性能实验表明，PM4Py预热后平均耗时随模型规模由0.0768秒增长到5.1961秒，且在72条以上消息流时增长更为明显。对于研究问题4，多场景人工案例组共有94/118个修复作用域和44/67个日志输入存在严格有效候选，说明方法能够在多数作用域中兼顾正常行为保留和异常行为接纳，但复合异常的联合修复仍是当前方法的主要限制；154个受控输入全部通过严格验证，证明了实现对目标异常类型和修复规则的覆盖能力。
